% =====================================================================
%  2.1 Likninger og formler — ekstra utfordring
% =====================================================================

\begin{infoboks}{Velg metode selv}
Arket bygger på 2.1 Likninger og formler. Vis regning og begrunnelse. Bruk eksakte svar.
Kontroller svaret ved å sette inn i den opprinnelige likningen.
\end{infoboks}

\nivaahode{2}

\oppg{U1}{2}{\uh}
Løs likningen og kontroller svaret.
\[ x - \frac{x-2}{3} = \frac{x+4}{2} - 1 \]

\oppg{U2}{2}{\uh}
Se på likningen $a(x-1) = 2x + 4$, der $a$ er en konstant.
\begin{tasks}(1)
  \task Løs likningen når $a = 3$.
  \task For hvilken verdi av $a$ har likningen ingen løsning?
  \task Finnes det en verdi av $a$ som gjør likningen til en identitet? Begrunn.
\end{tasks}

\oppg{U3}{2}{\uh}
Ved hvilken temperatur viser celsius og fahrenheit det samme tallet? Bruk $F = 1{,}8C + 32$.

\oppg{U4}{2}{\mh}
To taxiselskap tar ulik pris. Selskap A tar 49 kr i start og 16 kr per km. Selskap B tar 89 kr i start og 12 kr per km.
\begin{tasks}(1)
  \task Sett opp en formel for prisen hos hvert selskap etter $x$ km.
  \task Finn algebraisk hvor langt du må kjøre for at prisen blir lik.
  \task Tegn begge i GeoGebra. Hvilket selskap anbefaler du for en tur på 6 km og for en tur på 25 km?
\end{tasks}

\oppg{U5}{2}{\prog}
Programmet skal løse $ax + b = cx + d$ og skrive ut hvor mange løsninger det er. Fyll inn hullene.
\begin{kode}
def los(a, b, c, d):
    if a != c:
        x = (d - b) / (a - c)
        print("en losning:", x)
    elif ____:
        print("alle tall er losninger")
    else:
        print("____")
\end{kode}
\begin{tasks}(1)
  \task Hva må stå i de to hullene? Forklar hvorfor $a = c$ gir to ulike tilfeller.
  \task Hva skriver programmet ut for \lstinline[style=laereverk,numbers=none,frame=none]|los(3, 1, 1, 5)| og for \lstinline[style=laereverk,numbers=none,frame=none]|los(2, 1, 2, 4)|?
\end{tasks}

\nivaahode{3}

\oppg{U6}{3}{\uh}
\begin{tasks}(1)
  \task Vis at $\dfrac{1}{x} - \dfrac{1}{x+1} = \dfrac{1}{x(x+1)}$ for $x \neq 0$ og $x \neq -1$.
  \task Karl setter inn $x = 1$, $x = 2$ og $x = 3$, og alle gir lik verdi på begge sider. Forklar hvorfor det ikke er et bevis.
\end{tasks}

\oppg{U7}{3}{\uh}
Likningen $\sqrt{x^2} = x$ er sann for uendelig mange $x$, men den er ikke en identitet.
Forklar hvordan begge deler kan stemme. For hvilke $x$ er likningen sann?

\begin{oppsummering}
Én verdi som ikke stemmer, er nok til å vise at noe ikke er en identitet.
Å sette inn tall kan aldri bevise en identitet. Det krever algebra.
Ender du på $0 = 0$, er alle tall løsninger. Ender du på en usann påstand, er det ingen løsning.
\end{oppsummering}
